% CP-VI-0028
% target_chapter: VI/19 — Morphisms of Affine Schemes
% source_unit_id: IMP-DL-0B4CC237FA-P002
% source: Downloads/theory-of-algebraic-geometry-3.tex L302-L573
% solution_provenance: SOURCE_EXPLICIT_SOLUTION

\subsection*{Problem 12. The universal property of affine schemes}

Let \(A\) be a ring and let \((X,\mathcal O_X)\) be a scheme. Show that there is a bijection
\[
\alpha:
\operatorname{Hom}_{\Sch}(X,\operatorname{Spec} A)
\longrightarrow
\operatorname{Hom}_{\Rings}\bigl(A,\Gamma(X,\mathcal O_X)\bigr)
\]
given by associating to a morphism
\[
f:X\to \operatorname{Spec} A
\]
the induced map on global sections
\[
f^\#:
\Gamma(\operatorname{Spec} A,\mathcal O_{\operatorname{Spec} A})
\longrightarrow
\Gamma(X,\mathcal O_X).
\]

Since
\[
\Gamma(\operatorname{Spec} A,\mathcal O_{\operatorname{Spec} A})\cong A,
\]
this gives a ring homomorphism
\[
A\to \Gamma(X,\mathcal O_X).
\]

\textbf{Solution.}

We prove that the construction
\[
f\longmapsto f^\#
\]
gives a bijection
\[
\operatorname{Hom}_{\Sch}(X,\operatorname{Spec} A)
\cong
\operatorname{Hom}_{\Rings}(A,\Gamma(X,\mathcal O_X)).
\]

First, let
\[
f:X\to \operatorname{Spec} A
\]
be a morphism of schemes. A morphism of schemes consists of a continuous map of spaces together with a morphism of sheaves of rings
\[
f^\#:\mathcal O_{\operatorname{Spec} A}\to f_*\mathcal O_X.
\]
Taking global sections gives
\[
\Gamma(\operatorname{Spec} A,\mathcal O_{\operatorname{Spec} A})
\longrightarrow
\Gamma(\operatorname{Spec} A,f_*\mathcal O_X).
\]
But
\[
\Gamma(\operatorname{Spec} A,f_*\mathcal O_X)
=
\Gamma(X,\mathcal O_X).
\]
Also,
\[
\Gamma(\operatorname{Spec} A,\mathcal O_{\operatorname{Spec} A})\cong A.
\]
Hence every morphism
\[
f:X\to \operatorname{Spec} A
\]
determines a ring homomorphism
\[
\alpha(f):A\to \Gamma(X,\mathcal O_X).
\]
Thus \(\alpha\) is well-defined.

Now we construct the inverse map.

Suppose we are given a ring homomorphism
\[
\varphi:A\to \Gamma(X,\mathcal O_X).
\]
We must construct a morphism of schemes
\[
f_\varphi:X\to \operatorname{Spec} A.
\]

For every point \(x\in X\), we have natural maps
\[
\Gamma(X,\mathcal O_X)
\longrightarrow
\mathcal O_{X,x}
\longrightarrow
\kappa(x),
\]
where
\[
\kappa(x)=\mathcal O_{X,x}/\mathfrak m_x
\]
is the residue field of \(X\) at \(x\).

Composing these maps with \(\varphi\), we get a ring homomorphism
\[
A
\xrightarrow{\varphi}
\Gamma(X,\mathcal O_X)
\longrightarrow
\mathcal O_{X,x}
\longrightarrow
\kappa(x).
\]
Since \(\kappa(x)\) is a field, the kernel of this homomorphism is a prime ideal of \(A\). Therefore we define
\[
f_\varphi(x)
=
\ker\left(A\to \kappa(x)\right).
\]
Thus we obtain a set map
\[
f_\varphi:X\to \operatorname{Spec} A.
\]

We now prove that \(f_\varphi\) is continuous.

It is enough to compute the inverse image of a basic open subset
\[
D(a)\subseteq \operatorname{Spec} A.
\]
We have
\[
x\in f_\varphi^{-1}(D(a))
\]
if and only if
\[
a\notin f_\varphi(x).
\]
By definition of \(f_\varphi(x)\), this is equivalent to saying that the image of \(a\) in \(\kappa(x)\) is nonzero.

But the image of \(a\) in \(\kappa(x)\) is the image of the germ
\[
\varphi(a)_x\in \mathcal O_{X,x}
\]
in the residue field \(\kappa(x)\). Since \(\mathcal O_{X,x}\) is a local ring, an element is invertible if and only if its image in the residue field is nonzero. Hence
\[
a\notin f_\varphi(x)
\]
if and only if
\[
\varphi(a)_x\in \mathcal O_{X,x}^{\times}.
\]
Therefore
\[
f_\varphi^{-1}(D(a))
=
\{x\in X:\varphi(a)_x\in \mathcal O_{X,x}^{\times}\}.
\]
The locus where a section of a sheaf of rings is invertible is open. Hence
\[
f_\varphi^{-1}(D(a))
\]
is open in \(X\). Therefore \(f_\varphi\) is continuous.

Next we define the morphism of sheaves.

Let \(D(a)\subseteq \operatorname{Spec} A\) be a basic open subset. Then
\[
\mathcal O_{\operatorname{Spec} A}(D(a))=A_a.
\]
On the open subset
\[
f_\varphi^{-1}(D(a))\subseteq X,
\]
the section \(\varphi(a)\) is invertible. Hence the map
\[
\varphi:A\to \Gamma(X,\mathcal O_X)
\]
restricts to a map
\[
A\to \Gamma(f_\varphi^{-1}(D(a)),\mathcal O_X),
\]
and since \(\varphi(a)\) becomes invertible on this open set, it extends uniquely to a ring homomorphism
\[
A_a
\longrightarrow
\Gamma(f_\varphi^{-1}(D(a)),\mathcal O_X).
\]
Explicitly,
\[
\frac b{a^n}
\longmapsto
\frac{\varphi(b)}{\varphi(a)^n}
\]
on \(f_\varphi^{-1}(D(a))\).

These maps are compatible with restriction maps between basic opens. Therefore they glue to a morphism of sheaves
\[
f_\varphi^\#:\mathcal O_{\operatorname{Spec} A}\to (f_\varphi)_*\mathcal O_X.
\]
Thus we have constructed a morphism of ringed spaces
\[
f_\varphi:(X,\mathcal O_X)\to (\operatorname{Spec} A,\mathcal O_{\operatorname{Spec} A}).
\]

We must check that this is a morphism of locally ringed spaces. Let \(x\in X\), and put
\[
\mathfrak p=f_\varphi(x)\in \operatorname{Spec} A.
\]
The induced map on stalks is
\[
\mathcal O_{\operatorname{Spec} A,\mathfrak p}=A_{\mathfrak p}
\longrightarrow
\mathcal O_{X,x}.
\]
This map is local. Indeed, if \(a\in A\setminus \mathfrak p\), then the image of \(a\) in \(\kappa(x)\) is nonzero, so \(\varphi(a)_x\) is invertible in \(\mathcal O_{X,x}\). Therefore the localization map is well-defined. Moreover, the inverse image of the maximal ideal \(\mathfrak m_x\subseteq \mathcal O_{X,x}\) is precisely
\[
\mathfrak pA_{\mathfrak p}.
\]
Hence the stalk map
\[
A_{\mathfrak p}\to \mathcal O_{X,x}
\]
is a local homomorphism.

Therefore \(f_\varphi\) is a morphism of schemes.

Now we show that the two constructions are inverse to each other.

First start with a morphism of schemes
\[
f:X\to \operatorname{Spec} A.
\]
It induces a ring homomorphism
\[
A\to \Gamma(X,\mathcal O_X).
\]
Constructing the map on points from this homomorphism gives
\[
x\longmapsto
\ker\left(A\to \kappa(x)\right).
\]
But this is exactly the prime ideal corresponding to the image \(f(x)\in \operatorname{Spec} A\). Hence the reconstructed map agrees with \(f\) on points. The sheaf morphism also agrees with the original one, because on every basic open \(D(a)\), the morphism is forced by the corresponding localization map
\[
A_a\to \Gamma(f^{-1}(D(a)),\mathcal O_X).
\]
Therefore we recover the original morphism \(f\).

Conversely, start with a ring homomorphism
\[
\varphi:A\to \Gamma(X,\mathcal O_X).
\]
Construct \(f_\varphi:X\to \operatorname{Spec} A\), and then take the induced map on global sections. On the basic open \(\operatorname{Spec} A=D(1)\), the construction gives
\[
A=A_1
\longrightarrow
\Gamma(X,\mathcal O_X),
\]
and this map is exactly \(\varphi\). Hence we recover the original ring homomorphism.

Therefore the two constructions are inverse bijections.

Thus
\[
\boxed{
	\operatorname{Hom}_{\Sch}(X,\operatorname{Spec} A)
	\cong
	\operatorname{Hom}_{\Rings}\bigl(A,\Gamma(X,\mathcal O_X)\bigr)
}.
\]

This is the universal property of affine schemes.
